% !TeX root = ../../Book2.tex
% 纲要标题：### 9.4 外代数的应用
% 纲要位置：计划纲要.md，第 1086 行。

\section{外代数的应用}
\label{b2:sec:0904}

外代数既解释行列式，也提供记录子空间与构造复形的方法。本节先从子式推出伴随矩阵公式，并在任意交换环上证明 Cayley–Hamilton 定理；然后在域上用一个四维例子说明 Plücker 关系；最后给出 Koszul 微分的构造。三个应用分别用到最高外幂、二次外幂和整个分次外代数。

% 纲要要点（供目录核对）：
%
% * 余子式、伴随矩阵与 Cayley–Hamilton；系数递推和相邻项抵消
% * Plücker 坐标的初步例子；非零纯楔积恢复子空间、四维二次关系与特征二
% * Koszul 复形的外代数背景；收缩由分次乘法规则产生，反交换关系完整证明
% * 三变量自然关系与关系之间的相容性；复合为零不等于正合，自由消解的含义
% * 含单位序列的 Koszul 复形正合，构造及证明放在正文；正则序列理论归第三卷
%
% ---
%

\subsection{余子式、伴随矩阵与 Cayley–Hamilton 定理}
\label{b2:subsec:0904:01}

设 \(R\) 为交换环，\(A=(a_{ij})\in M_n(R)\)、\(n\ge1\)。删去第 \(i\) 行、第 \(j\) 列所得子矩阵记为 \(A^{\widehat i,\widehat j}\)。数
\[
C_{ij}=(-1)^{i+j}\det A^{\widehat i,\widehat j}
\]
称为对应的\term[daishu-yuzishi]{代数余子式}[cofactor]。当 \(n=1\) 时，删去后是空矩阵，行列式取 \(1\)。定义
\[
\operatorname{adj}(A)=(C_{ji})_{i,j}
\]
为 \(A\) 的伴随矩阵（adjugate matrix）\termidx[ban-sui-juzhen]{伴随矩阵}。这是代数余子式矩阵的转置，与第四章由首一多项式构造的友矩阵（companion matrix）是不同概念。
\symidx{\(\operatorname{adj}(A)\)：矩阵 \(A\) 的伴随矩阵}

\begin{proposition}[余子式展开与伴随矩阵]\label{b2:prop:adjugate-identity}
设 \(R\) 为交换环，\(A=(a_{ij})\in M_n(R)\)、\(n\ge1\)。令 \(C_{ij}=(-1)^{i+j}\det A^{\widehat i,\widehat j}\) 为代数余子式，\(\operatorname{adj}(A)=(C_{ji})_{i,j}\) 为伴随矩阵，其中 \(n=1\) 时空矩阵的行列式取 \(1\)。固定任意行 \(i\) 或列 \(j\)，有
\[
\det A=\sum_{k=1}^n a_{ik}C_{ik}
=\sum_{k=1}^n a_{kj}C_{kj}.
\]
进而
\begin{equation}\label{b2:eq:adjugate-identity}
A\operatorname{adj}(A)=\operatorname{adj}(A)A=(\det A)I_n.
\end{equation}
所以 \(A\) 可逆当且仅当 \(\det A\) 是 \(R\) 的单位，此时
\(A^{-1}=(\det A)^{-1}\operatorname{adj}(A)\)。
\end{proposition}
\begin{proof}
在\eqref{b2:eq:determinant-leibniz}中，按某一固定行所使用的列位置分组。含 \(a_{ik}\) 的项，其余因子的排列恰好遍历删去该行、该列后的置换。把第 \(i\) 行和第 \(k\) 列分别移到首位，共产生 \(i-1+k-1\) 个换位，故这一组的符号因子为 \((-1)^{i+k}\)。组内求和于是等于 \(a_{ik}C_{ik}\)，得到行展开。对转置矩阵使用行展开，再用\eqref{b2:eq:determinant-transpose}，得到列展开。

乘积 \(A\operatorname{adj}(A)\) 的 \((i,j)\) 项为 \(\sum_k a_{ik}C_{jk}\)。当 \(i=j\) 时是 \(\det A\)。当 \(i\ne j\) 时，我们仍要把这个和解释成一次行列式展开：把 \(A\) 的第 \(j\) 行替换为第 \(i\) 行，沿第 \(j\) 行展开时，代数余子式仍为 \(C_{jk}\)，行中的系数则变为 \(a_{ik}\)。所得行列式正是上述和，但矩阵有两行相同，故为零。这证明第一个乘积公式。第二个乘积公式通过替换第 \(i\) 列为第 \(j\) 列并作列展开得到。

若行列式为单位，式\eqref{b2:eq:adjugate-identity}给出双侧逆。若 \(A\) 可逆，则\cref{b2:thm:determinant-multiplicativity}给出 \(\det A\det A^{-1}=1\)，所以行列式为单位。
\end{proof}

下面通过比较多项式的系数重新证明 Cayley–Hamilton 定理\termidx[Cayley-Hamilton-dingli]{Cayley--Hamilton 定理}。\cref{b2:prop:operator-characteristic-minimal}利用域上的多项式环模分类，证明特征多项式是最小多项式的倍数，因而零化算子；这里直接使用行列式与伴随矩阵恒等式，把系数推广到任意交换环。

\begin{theorem}[Cayley–Hamilton]\label{b2:thm:exterior-cayley-hamilton}
若 \(R\) 为交换环、\(A\in M_n(R)\)、\(n\ge1\)，令
\[
\chi_A(t)=\det(tI_n-A)=t^n+c_{n-1}t^{n-1}+\cdots+c_0.
\]
则 \(\chi_A(A)=A^n+c_{n-1}A^{n-1}+\cdots+c_0I_n=0\)。
\end{theorem}
\begin{proof}
在交换环 \(R[t]\) 上应用\cref{b2:prop:adjugate-identity}，得到
\[
(tI_n-A)\operatorname{adj}(tI_n-A)=\chi_A(t)I_n.
\]
对标量多项式 \(\chi_A(t)\)，\(\chi_A(A)\) 当然有意义。但若要在左端按普通多项式代入的规则令 \(t=A\)，还须先核对矩阵系数与 \(A\) 交换：在矩阵系数多项式中，\(t\) 与所有系数交换，换成 \(A\) 后却不能未经证明地保留这一性质。因此这里保留每个矩阵因子的次序，改由系数关系消去伴随矩阵中的未知系数。

伴随矩阵的每个条目是 \(n-1\) 阶子式，所以次数至多 \(n-1\)。将它按 \(t\) 展开为
\[
\operatorname{adj}(tI_n-A)=B_0+B_1t+\cdots+B_{n-1}t^{n-1},
\qquad B_j\in M_n(R).
\]
比较 \(t\) 的各次系数，依次得到
\[
-AB_0=c_0I_n,\qquad
B_{j-1}-AB_j=c_jI_n\ (1\le j<n),\qquad B_{n-1}=I_n.
\]
这些关系把 \(c_jI_n\) 写成相邻两个 \(B_j\) 的组合，不必逐个求出 \(B_j\)。将中间第 \(j\) 个等式在左侧乘以 \(A^j\)，最后一式在左侧乘以 \(A^n\)，再连同首式求和，相邻的 \(B_j\) 项便有相同的左侧幂次而符号相反。右端为 \(\chi_A(A)\)，左端为
\[
-AB_0+\sum_{j=1}^{n-1}(A^jB_{j-1}-A^{j+1}B_j)
+A^nB_{n-1}=0.
\]
相邻项逐一抵消，故 \(\chi_A(A)=0\)。这一计算保持每个矩阵系数的位置，没有使用矩阵环中任意元素交换的假设。
\end{proof}

例如，对 \(A=\begin{pmatrix}a&b\\c&d\end{pmatrix}\)，有
\[
\chi_A(t)=t^2-(a+d)t+(ad-bc),
\qquad A^2-(a+d)A+(ad-bc)I_2=0.
\]
这些等式在 \(\mathbb Z\)、多项式环以及含幂零元的交换环上都成立。证明中用到的交换性属于标量环，而不是所有 \(2\times2\) 矩阵之间的交换性。

\subsection{Plücker 坐标}
\label{b2:subsec:0904:02}

本小节改设 \(k\) 为域、\(V=k^n\)。对 \(r\) 维子空间 \(W\subseteq V\)，选一组有序基 \(w_1,\ldots,w_r\)，写
\[
w_1\wedge\cdots\wedge w_r=\sum_{|I|=r}p_Ie_I.
\]
这些 \(p_I\) 是以 \(w_j\) 为列的矩阵的 \(r\) 阶子式。若换一组基，整个楔积乘以换基矩阵的非零行列式。因此，把所有 \(p_I\) 同时乘以同一个非零标量视为同一组坐标，就得到只依赖 \(W\) 的\term[Plucker-zuobiao]{Plücker 坐标}[Plücker coordinates]。这种只确定到共同非零倍数的坐标称为齐次坐标；子空间 \(W\) 确定的是 \(\bigwedge^rV\) 中的一条直线 \(k(w_1\wedge\cdots\wedge w_r)\)，而不是这条直线上的某个指定非零向量。

\ProseParagraphBreak
这些坐标确实能够辨认子空间。对非零纯楔积 \(\omega=w_1\wedge\cdots\wedge w_r\)，有
\begin{equation}\label{b2:eq:plucker-recovers-plane}
W=\{\,v\in V\mid v\wedge\omega=0\,\}.
\end{equation}
若 \(v\in W\)，按 \(w_j\) 展开即得楔积为零。若 \(v\notin W\)，则 \(v,w_1,\ldots,w_r\) 线性无关，\cref{b2:prop:wedge-linear-independence}说明它们的楔积非零。因此两个非零纯楔积只相差标量时，恢复出的子空间相同。这里的非零与纯楔积条件都不可省略；对一般外张量，右端不能无条件解释为一个 \(r\) 维子空间。\cref{b2:ex:09-plucker-plane-recovery}给出了由具体坐标恢复平面的计算。

\ProseParagraphBreak
还需判定哪些坐标真的来自一个子空间。例如，若一个 \(4\times2\) 矩阵的两列为 \(u,v\)，它的六个二阶子式同时来自这两列，因而必须满足彼此的相容关系。下面的二次式正是这种约束。

\begin{proposition}[四维空间中二重楔积的判据]\label{b2:prop:plucker-four-dimensional}
设 \(k\) 为任意域，\(e_1,\ldots,e_4\) 为 \(k^4\) 的标准基，并设
\[
\omega=\sum_{1\le i<j\le4}p_{ij}e_i\wedge e_j\in\bigwedge^2k^4.
\]
它可以写为 \(u\wedge v\)，当且仅当
\begin{equation}\label{b2:eq:plucker-quadric}
p_{12}p_{34}-p_{13}p_{24}+p_{14}p_{23}=0.
\end{equation}
此结论对任意特征的域成立。\(\omega\ne0\) 时，\(u,v\) 自动线性无关。
\end{proposition}
\begin{proof}
若 \(u=\sum u_ie_i\)、\(v=\sum v_ie_i\)，则 \(p_{ij}=u_iv_j-u_jv_i\)。代入\eqref{b2:eq:plucker-quadric}，得到
\[
\begin{aligned}
&(u_1v_2-u_2v_1)(u_3v_4-u_4v_3)\\
&\quad-(u_1v_3-u_3v_1)(u_2v_4-u_4v_2)\\
&\quad+(u_1v_4-u_4v_1)(u_2v_3-u_3v_2)=0.
\end{aligned}
\]
展开后的十二项按相同的 \(u_iu_jv_av_b\) 配对，符号相反，所以为零。

反过来，\(\omega=0\) 时取 \(u=0\)。若 \(\omega\ne0\)，至少一个 \(p_{ij}\ne0\)。重排基可使它成为 \(p_{12}\)：对相邻交换 \((1\ 2),(2\ 3),(3\ 4)\)，直接将下标重排并用 \(p_{ji}=-p_{ij}\)，可见关系左边都变为其负元，所以关系在重排后仍成立。令 \(a=p_{12}\ne0\)。我们先把待求的两列 \(u,v\) 的前两行选为 \((1,0)\)、\((0,a)\)，使 \(12\) 子式已经等于 \(a\)：
\[
\begin{pmatrix}
1&0\\
0&a\\
u_3&v_3\\
u_4&v_4
\end{pmatrix}.
\]
再匹配四个包含第 \(1\) 或第 \(2\) 行的子式，就得到
\[
v_3=p_{13},\qquad v_4=p_{14},\qquad
-a u_3=p_{23},\qquad -a u_4=p_{24}.
\]
前两行使这四个子式成为关于未知系数的线性方程。由于 \(a\ne0\)，它们唯一确定 \(u_3,u_4,v_3,v_4\)，即
\[
u=e_1-\frac{p_{23}}a e_3-\frac{p_{24}}a e_4,
\qquad v=a e_2+p_{13}e_3+p_{14}e_4.
\]
展开 \(u\wedge v\)，\(12,13,14,23,24\) 五个坐标分别就是给定的 \(p_{ij}\)，而 \(34\) 坐标为
\[
\frac{p_{13}p_{24}-p_{14}p_{23}}a=p_{34},
\]
最后一步恰用\eqref{b2:eq:plucker-quadric}：五个坐标已由所选两列实现，二次关系保证剩下的第六个坐标也匹配。故 \(u\wedge v=\omega\)。若 \(u,v\) 相关，其楔积为零，所以非零情形下它们独立。
\end{proof}

外幂中能写成一个纯楔积的元素称为\term[kefenjie-zhangliang]{可分解外张量}[decomposable exterior tensor]，也简称可分解张量。这里的“可分解”专指能写成一个楔积，区别于张量积中能写成一个纯张量。上面的关系表明，在 \(\bigwedge^2k^4\) 中，可分解外张量的坐标并不是任意六个标量。例如 \(e_1\wedge e_2+e_3\wedge e_4\) 的关系左端为 \(1\)，所以在任意域上都不可分解。一般 \(r\) 维子空间的 Plücker 坐标也满足一族二次关系；这里的 \(r=2,n=4\) 是只需一个非平凡二次方程的最小例子。

\ProseParagraphBreak
另一方面，简记 \(e_{ij}=e_i\wedge e_j\)、\(e_{1234}=e_1\wedge e_2\wedge e_3\wedge e_4\)。展开 \(\omega\wedge\omega\) 时，只有指标互不重复的交叉项可能非零。两个次数二的因子交换时产生的符号为 \((-1)^{2\cdot2}=1\)，所以 \(e_{12}\wedge e_{34}\) 与 \(e_{34}\wedge e_{12}\) 同号，都等于 \(e_{1234}\)。每一对这样的交叉项因而贡献两倍，得到
\[
\omega\wedge\omega
=2(p_{12}p_{34}-p_{13}p_{24}+p_{14}p_{23})
e_1\wedge e_2\wedge e_3\wedge e_4.
\]
特征不为二时，\(\omega\wedge\omega=0\) 与 Plücker 关系等价；特征二时，每个这样的平方都是零，包括刚才的不可分解张量。因而在特征二中仍应使用坐标关系本身来判定，不能把平方为零当作充分条件。

\subsection{Koszul 复形的外代数背景}
\label{b2:subsec:0904:03}

重新回到任意交换环 \(R\)。给定整数 \(n\ge1\) 与 \(a_1,\ldots,a_n\in R\)，以 \(e_1,\ldots,e_n\) 为 \(R^n\) 的标准基。生成元之间的关系由映射
\[
\lambda:R^n\longrightarrow R,\qquad e_i\longmapsto a_i
\]
的核记录：\(\sum_i r_ie_i\) 在核中，恰好表示 \(\sum_i r_i a_i=0\)。交换性立即给出关系 \(a_ie_j-a_je_i\)，因为它的像是 \(a_i a_j-a_j a_i=0\)。二次外幂中的 \(e_i\wedge e_j\) 正好记录这一对指标，更高外幂则用来组织这些关系之间的相容性。

\ProseParagraphBreak
对任意 \(R\) 模 \(M\)，为了把线性映射 \(\lambda:M\to R\) 延伸到整个外代数，我们要求它在一次分量上等于 \(\lambda\)，在标量上为零，并遵循分次 Leibniz 规则：作用于一个楔积时，可以作用于前一因子，也可以越过前一因子后作用于后一因子；越过次数 \(p\) 的因子时带上符号 \((-1)^p\)。次数为 \(-1\) 表示把 \(\bigwedge^pM\) 映到 \(\bigwedge^{p-1}M\)。若这样的算子连续作用两次总为零，即平方为零，就称为微分。本小节将所有负次外幂约定为零，使零次上的微分也有明确的目标。

\begin{proposition}[收缩微分]\label{b2:prop:exterior-contraction}
设 \(R\) 为交换含幺环，\(M\) 为幺性 \(R\) 模，\(\lambda:M\to R\) 为 \(R\) 线性映射，并约定 \(\bigwedge^pM=0\) 对 \(p<0\) 成立。存在唯一次数为 \(-1\) 的 \(R\) 线性分次导子
\[
\iota_\lambda:\bigwedge^pM\longrightarrow\bigwedge^{p-1}M,
\]
在 \(R=\bigwedge^0M\) 上取零，在 \(M=\bigwedge^1M\) 上等于 \(\lambda\)，且对 \(\xi\in\bigwedge^pM\)、\(\eta\in\bigwedge^qM\) 满足分次 Leibniz 规则
\[
\iota_\lambda(\xi\wedge\eta)
=\iota_\lambda(\xi)\wedge\eta+(-1)^p\xi\wedge\iota_\lambda(\eta).
\]
它在纯楔积上的公式为
\begin{equation}\label{b2:eq:contraction-formula}
\iota_\lambda(m_1\wedge\cdots\wedge m_p)
=\sum_{j=1}^p(-1)^{j-1}\lambda(m_j)
m_1\wedge\cdots\wedge\widehat{m_j}\wedge\cdots\wedge m_p.
\end{equation}
这里 \(p\ge1\)，帽号表示省去该因子。这个算子称为由 \(\lambda\) 确定的\term[shousuo]{收缩}[contraction]，并满足 \(\iota_\lambda^2=0\)。
\end{proposition}
\begin{proof}
若 \(D\) 是满足陈述中条件的分次导子，先把第一个因子分出，Leibniz 规则给出
\[
D(m_1\wedge\cdots\wedge m_p)
=\lambda(m_1)m_2\wedge\cdots\wedge m_p
-m_1\wedge D(m_2\wedge\cdots\wedge m_p).
\]
从 \(p=1\) 开始归纳，便强迫得到\eqref{b2:eq:contraction-formula}。纯楔积生成各次外幂，而零次上的作用已经指定，所以至多存在一个这样的算子。

反过来，式\eqref{b2:eq:contraction-formula}的右端对输入多线性。若 \(m_r=m_s\)、\(r<s\)，除删除第 \(r\) 或第 \(s\) 个因子的项以外，其余项仍有重复因子，均为零。删去第 \(r\) 个因子后，把原第 \(s\) 个因子移回第 \(r\) 个位置，共交换 \(s-r-1\) 次，所以该项的总符号是 \((-1)^{r-1+s-r-1}=(-1)^{s-2}\)，与删去第 \(s\) 个因子的项符号相反。二者系数中的 \(\lambda(m_r),\lambda(m_s)\) 相同，故抵消。这证明右端交替，外幂泛性质给出各次上的线性算子。把零次上的作用取为零，就得到次数为 \(-1\) 的 \(\iota_\lambda\)，其在一次分量上的作用是 \(\lambda\)。

当 \(p\) 或 \(q\) 为零时，所示乘法公式由 \(R\) 线性与标量上的零作用直接得到。对两个正次纯楔积的乘积，删除位置或者在前 \(p\) 项，或者在后 \(q\) 项。前者给出第一项，后者相对从 \(\eta\) 开始计数多出 \(p\)，产生 \((-1)^p\)，得到分次 Leibniz 规则。由双线性推广到所有齐次元素。

最后，两次收缩就是选两个不同位置，各施加一次 \(\lambda\) 并删去相应因子。每对位置有两种删除顺序，它们留下相同的楔积与相同的标量，只在符号上相反。具体说，固定 \(i<j\)，先删 \(i\) 再删原第 \(j\) 个因子的符号为 \((-1)^{i+j-3}\)，先删 \(j\) 再删 \(i\) 的符号为 \((-1)^{i+j-2}\)。两项的剩余因子次序相同，标量系数分别为 \(\lambda(m_i)\lambda(m_j)\) 及其反序，由交换性相等，故相互抵消。所有项成对消去，得到 \(\iota_\lambda^2=0\)。
\end{proof}
\symidx{\(\iota_\lambda\)：外代数上的收缩算子}

\begin{corollary}\label{b2:cor:contractions-anticommute}
设 \(R\) 为交换含幺环，\(M\) 为幺性 \(R\) 模，\(\lambda,\mu:M\rightarrow R\) 为 \(R\) 线性映射。由它们确定的外代数收缩算子满足
\[
\iota_\lambda\iota_\mu+\iota_\mu\iota_\lambda=0.
\]
这个反交换关系对任意特征成立。
\end{corollary}
\begin{proof}
由\eqref{b2:eq:contraction-formula}，\(\iota_{\lambda+\mu}\) 与 \(\iota_\lambda+\iota_\mu\) 在每个纯楔积上相同，在零次部分上也都为零，所以
\(\iota_{\lambda+\mu}=\iota_\lambda+\iota_\mu\)。对 \(\lambda+\mu\) 应用\cref{b2:prop:exterior-contraction}，得到
\[
0=\iota_{\lambda+\mu}^2
=\iota_\lambda^2+\iota_\lambda\iota_\mu
+\iota_\mu\iota_\lambda+\iota_\mu^2.
\]
两个平方项已为零，剩下的就是所求关系。这里只展开算子乘积并消去已知零项，没有除以二，因此特征二时也成立。
\end{proof}

令 \(n\ge1\)、\(M=R^n\)，选基 \(e_1,\ldots,e_n\)，并指定 \(\lambda(e_i)=a_i\in R\)。取 \(K_p=\bigwedge^pR^n\)、\(\mathrm{d}_p=\iota_\lambda\)，得到
\[
0\longrightarrow K_n\xrightarrow{\mathrm{d}_n}K_{n-1}\longrightarrow\cdots
\longrightarrow K_1\xrightarrow{\mathrm{d}_1}K_0\longrightarrow0.
\]
一列 \(R\) 模与 \(R\) 线性箭头若满足每两个相邻箭头的复合为零，就称为一个\term[lianfuxing]{链复形}[chain complex]；此处的条件 \(\mathrm{d}_{p-1}\mathrm{d}_p=0\) 已由\cref{b2:prop:exterior-contraction}证明。上面由 \(a_1,\ldots,a_n\in R\) 构造的链复形称为序列 \(a_1,\ldots,a_n\) 的\term[Koszul-fuxing]{Koszul 复形}[Koszul complex]，记为 \(K_\bullet(a_1,\ldots,a_n;R)\)。\symidx{\(K_\bullet(a_1,\ldots,a_n;R)\)：Koszul 复形}它在低次上的公式是
\[
\mathrm{d}_1(e_i)=a_i,\qquad \mathrm{d}_2(e_i\wedge e_j)=a_ie_j-a_je_i.
\]
\(K_1=R^n\) 记录这组理想生成元，\(K_2\) 记录由交换性自动产生的两两关系。因此 \(\operatorname{im}\mathrm{d}_1=(a_1,\ldots,a_n)\)，它的余核就是 \(R/(a_1,\ldots,a_n)\)；把这个商接到末端，是由 \(\mathrm{d}_1\) 的像决定的。得到增广复形
\[
0\longrightarrow K_n\longrightarrow\cdots\longrightarrow K_1
\xrightarrow{\mathrm{d}_1}R\longrightarrow R/(a_1,\ldots,a_n)\longrightarrow0.
\]
例如，当 \(n=2\) 时，未增广的复形为
\[
0\longrightarrow R\xrightarrow{\binom{-a_2}{a_1}}R^2
\xrightarrow{\begin{psmallmatrix}a_1&a_2\end{psmallmatrix}}R\longrightarrow0.
\]
相邻矩阵的乘积为 \(-a_1a_2+a_2a_1=0\)，\(\mathrm{d}_2\) 记录了最直接的两两关系。

\ProseParagraphBreak
当 \(n=3\)，写 \(a_1=a,a_2=b,a_3=c\)，并简记 \(e_{ij}=e_i\wedge e_j\)、\(e_{123}=e_1\wedge e_2\wedge e_3\)。三条两两关系为 \(ae_2-be_1\)、\(ae_3-ce_1\)、\(be_3-ce_2\)，而三次外幂给出
\[
\mathrm{d}_3(e_{123})=a e_{23}-b e_{13}+c e_{12}.
\]
它说明这三条关系本身也有相容关系，因为
\[
\begin{aligned}
\mathrm{d}_2\mathrm{d}_3(e_{123})
&=a(be_3-ce_2)-b(ae_3-ce_1)+c(ae_2-be_1)\\
&=0.
\end{aligned}
\]
于是三次外幂记录的是二次外幂所给关系之间的关系，更高次外幂继续以相同方式组织这些自然关系的相容性。

\ProseParagraphBreak
这些自然关系未必生成全部关系。取任意域 \(k\)，令 \(R=k[x]\)、\(a_1=a_2=x\)。此时 \(\mathrm{d}_1(u,v)=x(u+v)\)，\(\mathrm{d}_2(t)=(-xt,xt)\)。因为 \(R\) 是整环且 \(x\ne0\)，有
\[
\ker\mathrm{d}_1=R(-1,1),\qquad
\operatorname{im}\mathrm{d}_2=xR(-1,1)\subsetneq R(-1,1).
\]
严格包含来自 \(1\notin xR\)：关系 \((-1,1)\) 在核中，却不是 \(\mathrm{d}_2\) 的像。因此序列 \(x,x\) 的 Koszul 复形在 \(R^2\) 处不正合。

\ProseParagraphBreak
复合为零给出 \(\operatorname{im}\mathrm{d}_p\subseteq\ker\mathrm{d}_{p-1}\)；在相应位置正合则要求二者相等。增广复形在 \(R\) 及商模处正合，但在其他位置未必正合。一条增广链复形若在每个位置正合，且末端模之前的各项均为自由模，就称为这个末端模的\term[ziyou-xiaojie]{自由消解}[free resolution]。这里各 \(K_p=\bigwedge^pR^n\) 都是自由模，因此是否给出自由消解，关键在于尚待检验的正合性。

\ProseParagraphBreak
最短的例子是 \(n=1\)，复形就是 \(0\to R\xrightarrow{\times a_1}R\to0\)。若 \(k\) 为域，\(R=k[\varepsilon]/(\varepsilon^2)\)、\(a_1=\varepsilon\)，则左侧映射的核为非零理想 \((\varepsilon)\)，所以增广后也不是自由消解。哪些序列使增广后的 Koszul 复形成为 \(R/(a_1,\ldots,a_n)\) 的自由消解，将在第三卷第15.2—15.3节结合正则序列系统讨论。若一个生成元已经是单位，这里就能直接构造每个核元素的原像。

\begin{proposition}\label{b2:prop:koszul-unit-contractible}
设 \(R\) 为交换含幺环，\(n\ge1\)，\(a_1,\ldots,a_n\in R\)，且 \(a_1\) 是单位。令 \(e_1,\ldots,e_n\) 为 \(R^n\) 的标准基，\(K_\bullet=K_\bullet(a_1,\ldots,a_n;R)\) 为 Koszul 复形，其微分由 \(R\) 线性映射 \(\lambda:R^n\rightarrow R\)、\(\lambda(e_i)=a_i\) 的收缩给出。把 \(K_p=\bigwedge^pR^n\) 在 \(p<0\) 或 \(p>n\) 时取为零。定义次数为 \(+1\) 的 \(R\) 线性映射
\[
h_p:K_p\longrightarrow K_{p+1},
\qquad h_p(\xi)=a_1^{-1}e_1\wedge\xi.
\]
则在每个 \(0\le p\le n\) 上有
\[
\mathrm{d}_{p+1}h_p+h_{p-1}\mathrm{d}_p=\operatorname{id}_{K_p},
\]
其中边界处的映射取零。因此 \(K_\bullet\) 正合。
\end{proposition}
\begin{proof}
由\cref{b2:prop:exterior-contraction}的分次 Leibniz 规则及 \(\mathrm{d}_1(e_1)=a_1\)，对 \(\xi\in K_p\) 有
\[
\begin{aligned}
\mathrm{d}_{p+1}h_p(\xi)
&=a_1^{-1}\cdot\mathrm{d}_{p+1}(e_1\wedge\xi)\\
&=a_1^{-1}\bigl(a_1\xi-e_1\wedge\mathrm{d}_p(\xi)\bigr)\\
&=\xi-h_{p-1}\mathrm{d}_p(\xi).
\end{aligned}
\]
当 \(p=0\) 时，\(\mathrm{d}_0=0\)；当 \(p=n\) 时，\(e_1\wedge\xi\in\bigwedge^{n+1}R^n=0\)，上述等式在两个端点仍成立。这就证明所述恒等式，通常简写为 \(\mathrm{d}h+h\mathrm{d}=\operatorname{id}\)。

若 \(\mathrm{d}_p(\xi)=0\)，恒等式给出
\(\xi=\mathrm{d}_{p+1}(h_p(\xi))\)。这为每个闭元素（微分为零的元素）明确给出了成为边界（下一次微分的像）的原像。故 \(\ker\mathrm{d}_p\subseteq\operatorname{im}\mathrm{d}_{p+1}\)，反向包含由微分平方为零成立，复形在每处正合。
\end{proof}
